% !TeX program = xelatex
% 阿里云服务清单（报销附件）：UTF-8 编码，使用 XeLaTeX 编译。
% 修改下方基本信息和表格行，金额直接以“元”填写，并同步修改合计。
\documentclass[UTF8,a4paper,zihao=5,fontset=fandol]{ctexart}
\usepackage[margin=20mm]{geometry}
\usepackage{array,tabularx,longtable}
\usepackage{xcolor}
\definecolor{Accent}{HTML}{E85D04}
\definecolor{Muted}{HTML}{666666}
\setlength{\parindent}{0pt}
\setlength{\parskip}{7pt}
\setlength{\tabcolsep}{4pt}
\renewcommand{\arraystretch}{1.75}
\newcolumntype{L}[1]{>{\raggedright\arraybackslash}p{#1}}
\newcolumntype{R}[1]{>{\raggedleft\arraybackslash}p{#1}}
\newcolumntype{Y}{>{\raggedright\arraybackslash}X}
\pagestyle{plain}

% ---------- 基本信息 ----------
\newcommand{\Company}{合肥砰然科技有限公司}
\newcommand{\ExpensePeriod}{2026-09-09 至 2026-10-09}
\newcommand{\InvoiceNumber}{26342000002879875846}
\newcommand{\PreparedDate}{2026-09-09}
\newcommand{\Purpose}{为使用人本人的 ChatGPT 账户完成一次订阅周期的订阅购买及支付服务。}
\newcommand{\TotalAmount}{1470.00}

\begin{document}
\begin{center}
  {\sffamily\bfseries\LARGE 合肥砰然科技有限公司服务清单}\\[5pt]
  {\small\color{Muted}报销附件}
\end{center}
{\color{Accent}\rule{\textwidth}{1pt}}

\begin{tabularx}{\textwidth}{@{}L{20mm}Y L{20mm}Y@{}}
  单位名称 & \Company & 服务期间 & \ExpensePeriod \\
  发票号码 & \InvoiceNumber & 填制日期 & \PreparedDate \\
\end{tabularx}

\vspace{5pt}
\textbf{服务明细}\hfill{\small 金额单位：人民币元}

% ---------- 服务明细：复制或删除一行即可增删产品 ----------
% 六列依次为：产品名称、产品类型、订单状态、实付金额、订单创建时间、订单编号。
% 金额填本次报销所对应费用期间的实付金额；订单编号按实际记录填写。
\begingroup
\small
\setlength{\LTleft}{0pt}
\setlength{\LTright}{0pt}
\setlength{\LTpre}{4pt}
\setlength{\LTpost}{4pt}
\begin{longtable}{@{}L{25mm}L{25mm}L{17mm}R{20mm}L{25mm}L{\dimexpr\textwidth-112mm-40pt\relax}@{}}
  \hline
  \textbf{产品名称} & \textbf{产品类型} & \textbf{订单状态} & \textbf{实付金额} & \textbf{创建时间} & \textbf{订单编号} \\
  \hline
\endfirsthead
  \hline
  \textbf{产品名称} & \textbf{产品类型} & \textbf{订单状态} & \textbf{实付金额} & \textbf{创建时间} & \textbf{订单编号} \\
  \hline
\endhead
  ChatGPT Pro Coding Plan & AI 软件服务 & 已支付 & 1470.00 & 2026-09-09 & 5127727573274048846 \\
  \hline
  \multicolumn{3}{@{}l}{\textbf{合计}} & \textbf{\TotalAmount} & & \\
  \hline
\end{longtable}
\endgroup

\vspace{8pt}
\textbf{费用用途：}\Purpose

\end{document}
